\section{La capa de comunicación: cuándo hablar y cuándo reiniciar}
\label{sec:comunicacion}

Todos los consensos de este trabajo —el mapa de \cref{th:mapa}, los precios de
\eqref{eq:tatonnement}— se han supuesto periódicos, y \cref{le:muestras} tensa
las barreras con el periodo. Esa elección tiene un coste que conviene mirar de
frente, y una alternativa con garantías propias.

\subsection{Qué se paga por hablar a intervalos fijos}

\begin{observacion}[Dos objeciones al muestreo periódico]
\label{obs:objeciones-periodico}
La difusión periódica es fácil de implementar y de analizar, pero
\cite{zhao2024hybrid} señala dos inconvenientes. El primero es económico: se
transmite aunque no ocurra nada, de modo que en régimen estacionario —cuando los
precios ya han convergido y el mapa no cambia— se gasta ancho de banda y energía
en repetir información conocida. El segundo es de seguridad, y no aparece en
ninguna otra parte de este trabajo: la periodicidad hace el sistema
\emph{predecible}, y un adversario que conozca los instantes de transmisión
puede concentrar su interferencia exactamente en ellos. Un esquema aperiódico no
elimina esa amenaza, pero deja de regalar el calendario.
\end{observacion}

\subsection{Disparo por eventos con tiempo mínimo garantizado}

La alternativa es transmitir solo cuando la información retenida se ha quedado
obsoleta. La dificultad clásica es que un umbral puro admite comportamiento de
Zenón, y que las construcciones que lo evitan suelen dejar el tiempo mínimo
entre eventos sin cuantificar.

\begin{definicion}[Mecanismo de disparo con temporizador]
\label{def:etm}
Cada AMR mantiene el último valor difundido $\hat\chi_i$ y un temporizador
$\tau_i$ con $\dot\tau_i=1$. Difunde, poniendo $\hat\chi_i\leftarrow\chi_i$ y
$\tau_i\leftarrow0$, cuando
\[
  \lVert \chi_i-\hat\chi_i\rVert\ >\ \sigma\bigl(\lVert\chi_i\rVert+\varepsilon_0\bigr)
  \qquad\textbf{y}\qquad \tau_i\ \ge\ \tau^{\mathrm{MIET}} .
\]
\end{definicion}

\begin{proposicion}[Ausencia de Zenón por construcción]
\label{pr:miet}
Bajo \cref{def:etm} el intervalo entre dos difusiones consecutivas de un mismo
AMR es al menos $\tau^{\mathrm{MIET}}>0$. En consecuencia el número de eventos en
todo intervalo finito es finito, sin invocar términos de decaimiento
exponencial ni información de los vecinos en la condición de disparo.
\end{proposicion}

\begin{proof}
La segunda condición impide el disparo mientras $\tau_i<\tau^{\mathrm{MIET}}$, y
$\tau_i$ se anula solo al difundir; luego entre dos difusiones transcurre al
menos $\tau^{\mathrm{MIET}}$.
\end{proof}

\begin{observacion}[El temporizador ya estaba, cambia su papel]
\label{obs:temporizador-doble}
La variable que en el esquema periódico genera el instante de muestreo es la
misma que aquí impone el mínimo entre eventos \cite{zhao2024hybrid}. No se añade
maquinaria: se reinterpreta. Y con ello se atiende a la vez la exclusión de
Zenón y la limitación de \emph{hardware}, que impone un tiempo mínimo real entre
transmisiones y que un análisis puramente asintótico suele pasar por alto. Es la
misma lógica de \cref{th:dwell} —un tiempo mínimo de permanencia— aplicada ahora
a la comunicación en lugar de a la topología.
\end{observacion}

\begin{observacion}[Cuánto se ahorra y qué se paga]
\label{obs:etm-ahorro}
Sobre el consenso de precios en un anillo de ocho AMR durante ocho segundos, la
difusión periódica a $20$~ms consume $3200$ mensajes y deja una dispersión final
de $6{,}5\times10^{-3}$. Con $\sigma=0{,}05$ y $\tau^{\mathrm{MIET}}=100$~ms se
consumen $90$ mensajes —un $97\%$ menos— con dispersión final
$4{,}8\times10^{-2}$.

El ahorro es grande y el precio también: la dispersión residual empeora en un
factor siete, y es proporcional a $\sigma$. Conviene por tanto no presentar el
disparo por eventos como una mejora sin contrapartida. Es un intercambio entre
ancho de banda y precisión de acuerdo, y el punto de operación se elige sabiendo
que el error residual no tiende a cero mientras $\sigma>0$.
\end{observacion}

\begin{observacion}[Lo que el tiempo mínimo \emph{no} acota]
\label{obs:miet-no-acota-max}
Aquí hay una trampa que conviene declarar. $\tau^{\mathrm{MIET}}$ acota el
intervalo \emph{mínimo}, que es lo que exige el \emph{hardware}. La seguridad
exige lo contrario: \cref{le:muestras} tensa las barreras con el intervalo
\emph{máximo}, y el disparo por eventos no lo acota. En el experimento anterior
el intervalo más largo observado fue de $3{,}15$~s, ciento cincuenta veces el
periodo del esquema periódico.

La consecuencia es de diseño y es obligatoria: la capa de seguridad no puede
depender de información difundida por eventos sin un \emph{plazo máximo de
transmisión} añadido —una difusión forzada cuando $\tau_i$ alcanza
$\tau^{\max}$— y es $\tau^{\max}$, no $\tau^{\mathrm{MIET}}$, el que entra en el
tensado de \cref{le:muestras}. Con esa salvaguarda el esquema tiene dos
temporizadores, uno por abajo y otro por arriba, y ambos son necesarios por
razones distintas.
\end{observacion}

\subsection{Un solo AMR que informe mal se lleva el consenso entero}

\Cref{obs:objeciones-periodico} señala una vulnerabilidad de calendario. Hay otra
más elemental: todos los protocolos de
consenso empleados aquí suponen que cada AMR difunde su valor \emph{verdadero}.
Un nodo averiado —o comprometido— que difunda un valor arbitrario no está
contemplado en ninguna hipótesis.

\begin{observacion}[La media no resiste a un solo desviado]
\label{obs:consenso-resiliente}
Sobre un anillo de ocho AMR con dos vecinos por lado y valores iniciales en
$[0,10]$, un único nodo que difunda de forma persistente el valor $50$ arrastra a
los \emph{siete} restantes exactamente a $50$. No a un valor sesgado: al valor
que él elija. El promedio ponderado no tiene defensa alguna, porque cada nodo
leal trata el dato falso como cualquier otro.

Basta sustituir la media por una regla que descarte los $F$ valores extremos por
cada lado antes de promediar para que los leales converjan a $6{,}775$, dentro
del rango de sus valores iniciales. La condición para poder aplicarla es tener
más de $2F$ vecinos, lo que impone una densidad mínima al grafo de comunicación
\cite{mahmoud2021advanced}.
\end{observacion}

\begin{observacion}[Qué se pierde al blindarlo, y dónde duele]
\label{obs:resiliente-precio}
El blindaje no es gratuito y su coste es preciso: la regla filtrada \emph{ya no
converge a la media}. Converge a un valor dentro del rango de los leales, que es
lo que se puede garantizar sin saber quién miente, pero no al promedio que
\cref{th:mapa} necesita para igualar a la fusión bayesiana. Es decir, blindar la
cartografía cuesta exactitud en el mismo sentido en que
\cref{obs:etm-ahorro} la cuesta por ancho de banda.

Y hay una consecuencia sobre el diseño del grafo que tira en contra de
\cref{obs:esparsificar}: exigir más de $2F$ vecinos impide adelgazar hasta el
grafo de Gabriel. Robustez ante nodos que mienten y economía de aristas son
objetivos opuestos y no se optimizan a la vez.

Se declara como limitación, no como resultado: los protocolos de este trabajo
suponen difusión veraz, y esa hipótesis debe añadirse a las que
\cref{obs:complejidad-propia} enumera. El remedio existe y es conocido; su
integración con la dinámica de multiplicadores de \cref{s:equilibrado} —que
además exige equilibrio de pesos— no es inmediata y no se aborda aquí.
\end{observacion}

\subsection{Contabilidad: qué cuesta cada tarea}

Los números de \cref{obs:etm-ahorro} son medidas de un experimento. Conviene
darles el marco formal que existe para ello, porque distingue tres cosas que se
confunden con facilidad \cite{bullo2009distributed}.

\begin{definicion}[Complejidad temporal y de comunicación]
\label{def:complejidad}
Sea $\Tcal$ una tarea de coordinación y $\Ccal\Ccal$ una ley de control y
comunicación compatible con ella. Se definen:
\begin{enumerate}[label=(\roman*),leftmargin=2.2em,topsep=0.2em,itemsep=0.15em]
  \item la complejidad temporal desde una condición inicial,
  $\mathrm{TC}(\Tcal,\Ccal\Ccal,x_0)=\inf\{\ell\mid
  \Tcal(x(k))=\text{cierto},\ \forall k\ge\ell\}$: la ronda a partir de la cual
  la tarea se cumple \emph{y se sigue cumpliendo};
  \item su peor caso $\mathrm{TC}(\Tcal,\Ccal\Ccal)$, tomando supremo sobre las
  condiciones iniciales admisibles;
  \item la complejidad de la \emph{tarea},
  $\mathrm{TC}(\Tcal)=\inf\{\mathrm{TC}(\Tcal,\Ccal\Ccal)\mid \Ccal\Ccal
  \text{ compatible}\}$, ínfimo sobre todas las leyes admisibles.
\end{enumerate}
Análogamente, la complejidad de comunicación total acumula el número de mensajes
hasta $\mathrm{TC}$, y la media lo divide por ese horizonte. Para tareas que
deben \emph{mantenerse}, la magnitud pertinente es la media en horizonte
infinito, $\lim_{\tau\to\infty}\tau^{-1}\sum_{\ell=0}^{\tau}|M(x(\ell))|$.
\end{definicion}

\begin{observacion}[El tercer nivel es el que permite afirmar optimalidad]
\label{obs:complejidad-tarea}
Los dos primeros niveles miden un algoritmo; el tercero mide el \emph{problema}.
Solo con él puede decirse que una construcción es esencialmente insuperable, y
por eso es el que echaba en falta \cref{obs:falta-optimalidad}: la cota que allí
se declara abierta es, en este vocabulario, una cota inferior sobre
$\mathrm{TC}(\Tcal)$ y sobre la complejidad de comunicación de la tarea, no sobre
la de una política concreta.
\end{observacion}

\begin{observacion}[Qué se puede afirmar hoy y qué no]
\label{obs:complejidad-propia}
Aplicado a este trabajo, el balance es desigual y conviene exhibirlo.

Del consenso de mapa y de precios se conoce la complejidad temporal en el
sentido (ii): la tasa de \cref{th:conmutada} la acota en función del diámetro y
de la conectividad algebraica, y \cref{th:jerarquia-grafos} dice cómo se degrada
al adelgazar el grafo. Su complejidad de comunicación en horizonte infinito es
exactamente lo que mide \cref{obs:etm-ahorro}, y el disparo por eventos la
reduce en dos órdenes de magnitud a cambio de precisión.

Del reclutamiento se conoce menos. \Cref{co:terminacion} garantiza terminación,
que es finitud de $\mathrm{TC}$, pero no da cota: el número de revisiones
depende del catálogo y de las cuencas del potencial, y \cref{obs:gibbs-precio}
muestra que el tiempo de escape crece exponencialmente con $\beta$. No hay, por
tanto, una cota de peor caso, y menos aún una cota de tarea.

Declararlo así separa lo medido de lo demostrado. Los números de este trabajo
son complejidades de comunicación de \emph{instancias} concretas; ninguna es
todavía una cota sobre el problema.
\end{observacion}

\subsection{Convergencia en tiempo finito, y por qué no se adopta}

\Cref{obs:complejidad-propia} declara que no hay cota de tiempo para el
reclutamiento. Existe una técnica que la daría, y conviene examinarla en lugar de
ignorarla.

\begin{lema}[Cota de tiempo de asentamiento]
\label{le:tiempo-finito}
Sea $V$ definida positiva y de clase $C^1$. Si existen $c>0$ y
$\alpha\in(0,1)$ con $\dot V+cV^{\alpha}\le0$, entonces $V$ alcanza el origen en
tiempo finito \cite{wang2024adaptive}, acotado por
\[
  T\ \le\ \frac{V^{1-\alpha}(x_0)}{c\,(1-\alpha)} .
\]
\end{lema}

\begin{observacion}[La cota es ajustada]
\label{obs:tiempo-finito-num}
Sobre $\dot x=-cx^{\alpha}$ la cota se alcanza con igualdad: con $c=1$,
$\alpha=0{,}5$, $x_0=4$ se mide $T=4{,}0000$ frente a una cota de $4{,}0000$; con
$c=2$, $x_0=9$, se mide $3{,}0000$ frente a $3{,}0000$. No es una estimación
conservadora sino el tiempo exacto.

Llevado al consenso, sustituir el acoplamiento lineal por la potencia fraccionaria
$-\sum_j a_{ij}\operatorname{sign}(x_i-x_j)\lvert x_i-x_j\rvert^{\alpha}$ hace que
la dispersión se anule \emph{exactamente}: con $\alpha=0{,}5$ en $2{,}89$~s y con
$\alpha=0{,}7$ en $4{,}09$~s, mientras que el protocolo lineal no la alcanza en
sesenta segundos. Eso es justo lo que falta en \cref{obs:complejidad-propia}.
\end{observacion}

\begin{observacion}[Por qué, aun así, no se adopta aquí]
\label{obs:tiempo-finito-precio}
La potencia fraccionaria no es Lipschitz en el origen, y esa es exactamente la
propiedad que da el tiempo finito. Bajo muestreo se paga: con paso $h$, la
dispersión residual medida es $2\times10^{-8}$ para $h=10^{-4}$,
$2\times10^{-6}$ para $10^{-3}$, $2\times10^{-4}$ para $10^{-2}$ y
$5\times10^{-3}$ para $5\times10^{-2}$ —crece como $h^{2}$— frente a
$\sim3\times10^{-9}$ del protocolo lineal, \emph{peor en los cuatro pasos
ensayados}. Y con $\alpha=0{,}3$ el esquema no converge en absoluto dentro del
horizonte.

La lectura es la misma que en \cref{obs:ascenso-muestreo} y en
\cref{le:muestras}: \emph{la garantía es de tiempo continuo y no
sobrevive a la discretización sin condiciones}. Aquí no se degrada suavemente
—se convierte en un entorno de castañeteo—, y en un sistema que muestrea a
$20$~ms el remedio sería peor que la enfermedad. Se declara como técnica
disponible y descartada, con el motivo, en lugar de adoptarla por la elegancia
de su enunciado.
\end{observacion}

\subsection{Reiniciar el ajuste de precios}

El ajuste \eqref{eq:tatonnement} es integral, y un integrador aporta retardo de
fase. La teoría de control por reinicio muestra que reiniciar el estado del
controlador en instantes escogidos reduce ese retardo sin alterar la respuesta
en amplitud, con lo que mejora el transitorio sin renunciar a la estabilidad
asintótica \cite{zhao2024hybrid}.

\begin{proposicion}[Reinicio del acumulador de precio]
\label{pr:reset}
Sea el ajuste con acción integral sobre el exceso de demanda. Si el acumulador
se pone a cero cuando el exceso cambia de signo, el punto de reposo no se altera
—porque en él el exceso es nulo y el acumulador no contribuye— mientras que el
sobreimpulso del transitorio se reduce.
\end{proposicion}

\begin{observacion}[Medida del efecto]
\label{obs:reset-num}
Sobre una planta de precios de segundo orden con controlador proporcional
integral, el sobreimpulso pasa de $46{,}1\%$ a $24{,}1\%$ al añadir el reinicio,
una reducción del $48\%$. A los doce segundos el lazo lineal sigue a $1{,}16$
del valor de referencia mientras el lazo con reinicio está en $1{,}01$. No se
afirma mejora del tiempo de establecimiento: en este ensayo la métrica del $2\%$
satura en el horizonte para ambos, y declararla mejorada sería leer el
experimento con más generosidad de la que permite.
\end{observacion}

\begin{observacion}[Hay una alternativa continua, y es mejor]
\label{obs:anticipacion}
El reinicio no es la única forma de corregir el transitorio, ni la más
conveniente. La teoría de dinámicas de pago ofrece el modelo con
\emph{suavizado y anticipación} \cite{martinezPiazuelo2025gne}:
\[
  \dot q=\tau\bigl(f(x)-q\bigr),
  \qquad
  p=\alpha_0 f(x)+\alpha_1 q+\alpha_2\dot q,
  \qquad \alpha_0+\alpha_1=1,\ \alpha_2\ge0,
\]
donde $q$ es una versión filtrada del pago y el término $\alpha_2\dot q$ anticipa
hacia dónde se dirige. Satisface por construcción tres de los cuatro requisitos
de la condición de estabilidad, y el cuarto —la $\delta$-antipasividad— bajo
condiciones suficientes explícitas.

Medido sobre el mismo lazo de precios de \cref{obs:reset-num}, el sobreimpulso
pasa de $46{,}1\%$ sin corrección a $35{,}1\%$ con $\alpha_2=0{,}10$,
$22{,}5\%$ con $0{,}25$ y $9{,}4\%$ con $0{,}50$, alcanzando a los doce segundos
el valor exacto de referencia. La anticipación bate al reinicio —$9{,}4\%$ frente
a $24{,}1\%$— y lo hace con un mecanismo continuo, no híbrido, y con garantías
dentro del mismo marco disipativo que sostiene \cref{th:smallgain}.

Y una advertencia obtenida al probarlo: \emph{combinar} ambos empeora las cosas.
Con anticipación $0{,}25$ más reinicio, el sobreimpulso baja a $11{,}6\%$ pero el
lazo termina en $0{,}679$ en lugar de en $1$: las dos correcciones se suman y
sobrecorrigen. Debe elegirse una.
\end{observacion}

\begin{observacion}[Dónde reiniciar en esta arquitectura]
\label{obs:reset-donde}
El instante natural de reinicio no es solo el cruce por cero del exceso. Cada
transición del autómata de \cref{def:automata} invalida el acumulado: al formarse
una coalición cambia el mapa de agarre, y el precio acumulado bajo la
composición anterior deja de corresponder a la demanda presente. Reiniciar en la
guarda evita arrastrar un integral obsoleto al nuevo régimen. Es una decisión de
implementación, no un teorema, y así se declara.
\end{observacion}

\subsection{Quién conoce la demanda}

Hay una suposición que este trabajo ha venido usando sin nombrarla: que la
demanda de \emph{wrench} $\Wdem_\ell$, generada por la trayectoria deseada de la
carga, está disponible para todos los miembros de la coalición. No lo está. La
referencia se origina en un punto —el AMR que aceptó la tarea, o un enlace con
el planificador— y los demás no tienen por qué verla.

\begin{definicion}[Observador distribuido de la referencia]
\label{def:observador-dist}
Sea la referencia generada por $\dot v_0=S_0v_0$ con $S_0$ conocida. Cada AMR
mantiene una estimación $v_i$ que evoluciona según
\[
  \dot v_i\;=\;S_0v_i\;+\;L_0\sum_{j\in\Ncal_i\cup\{0\}}a_{ij}\,W_0\,(v_j-v_i),
\]
es decir, una copia local de la dinámica de la referencia corregida por consenso
con los vecinos, donde $a_{i0}\neq0$ solo para los AMR con acceso directo
\cite{cai2022cooperative}.
\end{definicion}

\begin{proposicion}[Convergencia bajo la conectividad que ya se exige]
\label{pr:observador-dist}
El error $\tilde v=\operatorname{col}(v_i-v_0)$ obedece
$\dot{\tilde v}=(I\otimes S_0-H_\sigma\otimes L_0W_0)\tilde v$, con
$H_\sigma=L_\sigma+\operatorname{diag}(a_{10},\dots,a_{N0})$. Existe $L_0$ que lo
hace exponencialmente estable si y solo si $H_\sigma$ tiene todos sus autovalores
con parte real positiva, lo que equivale a que el grafo contenga un árbol
generador con raíz en el origen de la referencia.
\end{proposicion}

\begin{observacion}[No hace falta ninguna hipótesis nueva]
\label{obs:observador-encaja}
La condición de \cref{pr:observador-dist} es exactamente la de
\cref{th:conmutada}, que ya se exige para el consenso de mapa y de precios. Es
decir: la propagación de la referencia se obtiene gratis sobre la conectividad
que la arquitectura ya paga.

Comprobación sobre una cadena de cuatro AMR con la referencia oscilante: con una
sola arista al origen, $\min\operatorname{Re}\lambda(H)=1{,}0$ y el error final es
$8{,}5\times10^{-4}$; sin ninguna arista, $\min\operatorname{Re}\lambda(H)=0$ y el
error se estanca en $1{,}46$. El resultado no es que convenga tener el árbol: es
que sin él la estimación no converge en absoluto, por bien diseñado que esté
$L_0$.
\end{observacion}

\begin{observacion}[La fuga del ajuste de precios es un modelo interno ausente]
\label{obs:fuga-modelo-interno}
El principio del modelo interno dice que un esquema prealimentado no logra
seguimiento exacto cuando los parámetros son inciertos, mientras que un
controlador que contenga una copia de la dinámica de la señal a seguir —y que se
alimente del \emph{error}— conserva la regulación bajo perturbaciones
paramétricas \cite{cai2022cooperative}. La comprobación es contundente: sobre
una planta de primer orden con referencia constante, un error del $40\%$ en el
parámetro deja al prealimentado con $28{,}6\%$ de error permanente, y del $-40\%$
con $66{,}7\%$; el modelo interno —un integrador del error— deja
$10^{-14}$ en los tres casos.

Esto tiene una lectura incómoda y útil sobre el propio ajuste
\eqref{eq:tatonnement}. Su término de fuga $d_p>0$ hace que el reposo sea
$p_\ell=\alpha(G\bm\lambda-\Wdem_\ell)/d_p$, de modo que el exceso de demanda
\emph{no} se anula: queda un residuo proporcional a $d_p$. La fuga es
precisamente lo que impide que el ajuste sea un modelo interno de una demanda
constante. Hay, por tanto, un compromiso que conviene declarar: la fuga aporta amortiguamiento y robustez a la red de precios
—\cref{obs:alcance-poblacional} depende de ella— y cuesta exactitud en el
vaciado del mercado. Reducir $d_p$ acerca el vaciado exacto y aproxima la red de
precios a su frontera de estabilidad. Es un parámetro de diseño con dos efectos
opuestos, y conviene fijarlo sabiéndolo.
\end{observacion}

\subsection{Regular una salida común no es igualar motores}

\Cref{def:observador-dist} propaga la referencia; falta decir qué hace cada AMR
con ella. La tentación de la literatura de consenso es igualar estados. Con AMR
heterogéneos eso es un error: lo que debe coincidir es la \emph{salida} —el
\emph{twist} de la carga que cada uno imprime—, no los estados ni los pares
\cite{mayorga2026m1}.

\begin{proposicion}[Error de regulación con referencia local]
\label{pr:regulacion-local}
Sea el AMR $i$ con modelo local $\dot x_i=A_ix_i+B_iu_i+E_i\zeta$,
$e_i=C_ix_i+D_iu_i+F_i\zeta$, y la referencia común $\dot\zeta=S_0\zeta$.
Supóngase que las ecuaciones reguladoras
\[
  \Pi_iS_0=A_i\Pi_i+B_i\Gamma_i+E_i,\qquad C_i\Pi_i+D_i\Gamma_i+F_i=0
\]
tienen solución y que $A_i+B_iK_i$ es Hurwitz con
$\lVert e^{(A_i+B_iK_i)t}\rVert\le M_ie^{-\lambda_it}$. Con la ley
$u_i=K_i(x_i-\Pi_i\hat\zeta_i)+\Gamma_i\hat\zeta_i$, y llamando
$\tilde x_i=x_i-\Pi_i\zeta$, $\tilde\zeta_i=\hat\zeta_i-\zeta$,
\[
  \dot{\tilde x}_i=(A_i+B_iK_i)\tilde x_i+B_i(\Gamma_i-K_i\Pi_i)\tilde\zeta_i,
\]
\[
  \lVert\tilde x_i(t)\rVert\le M_ie^{-\lambda_it}\lVert\tilde x_i(0)\rVert
  +M_i\lVert B_i(\Gamma_i-K_i\Pi_i)\rVert\int_0^{t}e^{-\lambda_i(t-s)}\lVert\tilde\zeta_i(s)\rVert\,ds .
\]
\end{proposicion}

\begin{observacion}[Comprobado con dos AMR distintos]
\label{obs:regulacion-num}
Sobre dos AMR con dinámicas lineales diferentes siguiendo el mismo oscilador de
referencia, las soluciones reguladoras difieren —$\Gamma_0=(3{,}19;\,1{,}34)$
frente a $\Gamma_1=(0{,}89;\,0{,}28)$— y aun así ambos llevan el error de salida
a $2{,}7\times10^{-5}$ con un error de observador que decae. La cota es
esencialmente ajustada: la razón entre el error medido y la cota alcanza $1{,}00$
y $1{,}01$ tomando $M_i=1$, lo que indica que en sistemas no normales $M_i$ debe
tomarse ligeramente mayor que uno.

La lectura para la arquitectura es doble. Primera, el término que excita al
error es el del \emph{observador}, $\tilde\zeta_i$: la calidad de
\cref{pr:observador-dist} entra directamente en la regulación, y el
condicionamiento de \cref{th:saltos} fija cuántos saltos hacen falta para que
ese término sea pequeño. Segunda, y es la que corrige una lectura fácil de
\cref{obs:variable-coordinacion}: la variable de coordinación es el precio y la
referencia, y \emph{nada más}. Estados, pares y referencias internas $\Pi_i\zeta$
son distintos por diseño entre AMR heterogéneos, y forzar su consenso sería
imponer una restricción que la regulación no necesita y que la heterogeneidad
no admite. La generalización distribuida exige observadores, hipótesis de red y
regulabilidad; no equivale a tratar como independientes cuerpos unidos por
contacto.
\end{observacion}

\subsection{Una referencia que también escucha al equipo}

\Cref{def:observador-dist} propaga la referencia hacia los AMR en un solo
sentido: la referencia avanza y los seguidores la persiguen. Con contacto y
batería eso es un error de diseño: una referencia que sigue avanzando mientras
el equipo pierde contacto no es una variable de cooperación sino una orden
temporal. La corrección es hacer que el reloj de progreso escuche al error
colectivo, $ds/dt=1/(1+z_f)$ con $z_f\ge0$ creciente con ese error
\cite{ren2008distributed,mayorga2026m1}.

\begin{proposicion}[Reloj seguro para una familia reducida de errores]
\label{pr:reloj-seguro}
Supóngase $\dot e_i=v-k_ie_i$, $k_i>0$, con $0\le e_i\le\epsilon_i$, y progreso
$\dot s=v$. El mando
\[
  v=\min\Bigl\{v_{\mathrm{nom}},\ \min_i\bigl[k_ie_i+\alpha_i(\epsilon_i-e_i)\bigr]\Bigr\},
  \qquad \alpha_i>0,\ v_{\mathrm{nom}}>0,
\]
preserva el intervalo de errores; además
$v\ge v_{\min}:=\min\{v_{\mathrm{nom}},\min_i\epsilon_i\min(k_i,\alpha_i)\}>0$,
y una longitud $S$ se completa en no más de $S/v_{\min}$.
\end{proposicion}

\begin{proof}
En $e_i=0$, $\dot e_i=v\ge0$; en $e_i=\epsilon_i$ la restricción impone
$v\le k_i\epsilon_i$, luego $\dot e_i\le0$. Cada término de la cota es afín en
$e_i$ y su mínimo en el intervalo es $\epsilon_i\min(k_i,\alpha_i)$; integrar
$\dot s\ge v_{\min}$ da el tiempo.
\end{proof}

\begin{observacion}[Comprobado, y con su límite físico]
\label{obs:reloj-num}
Con $k=(2;\,0{,}5)$, $\alpha=(2;\,2)$, $\epsilon=(0{,}2;\,0{,}2)$ y
$v_{\mathrm{nom}}=0{,}3$: $v_{\min}=0{,}1$, un metro se recorre en $9{,}93$~s
frente a la cota de $10$~s, la velocidad observada nunca baja de $0{,}1000$ y el
error máximo toca exactamente $0{,}2000$ sin cruzarlo. Con reloj fijo
$v=v_{\mathrm{nom}}$ el metro se recorre en $3{,}33$~s, pero el error alcanza
$0{,}52$: dos veces y media el límite de seguimiento.

El traslado a la arquitectura es que el progreso a lo largo de la ruta pasa a ser
una magnitud negociada junto al esfuerzo y a los precios, sujeta a la capacidad
de seguimiento, en lugar de un reloj externo. Y el límite del resultado es
preciso: son dos ecuaciones de error de primer orden, no AMR ni cargas; una
pared o un contacto perdido pueden imponer $v=0$, y entonces desaparece la cota
de progreso. Trasladarla a la planta exige derivar una envolvente de
seguimiento válida para ella, que es lo que \cref{def:testigo} llama tubo.
\end{observacion}

\begin{observacion}[El bloque antisimétrico del primal--dual]
\label{obs:bloque-antisimetrico}
Una precisión afín sobre el cálculo de precios. Con coste convexo $\Phi$ y
balance $Au=b$, el operador de silla
$F(u,\lambda)=(\nabla\Phi(u)+A^{\!\top}\lambda,\ b-Au)$ contiene el bloque
antisimétrico $\bigl(\begin{smallmatrix}0&A^{\!\top}\\-A&0\end{smallmatrix}\bigr)$:
aunque $\Phi$ tenga potencial, iterar $F$ no es descender una función en todas
sus coordenadas, y el método de Euler diverge ya en $L(x,\lambda)=x\lambda$, cuyo
operador es una rotación pura. Es la parte \emph{rotacional} del juego, y es la
razón de fondo por la que el ajuste de precios necesita un paso de predicción y
otro de corrección, como el extragradiente ya adoptado en
\cref{sec:informacion2} \cite{mayorga2026m1}. Los resultados extraproximales de
la literatura se formulan para su clase de juegos y regularizaciones; no se
trasplantan a todo contacto no convexo.
\end{observacion}

\subsection{Lo que este marco ofrece y no se usa}

\begin{observacion}[El consenso bipartito no aplica]
\label{obs:bipartito}
Buena parte de \cite{zhao2024hybrid} desarrolla el \emph{consenso bipartito}
sobre grafos con signo, donde algunas relaciones son antagónicas y los agentes
convergen a la misma magnitud con signos opuestos, bajo la condición de que el
grafo sea estructuralmente equilibrado. La tentación de aplicarlo aquí es que
los AMR compiten por las cargas. Debe resistirse: la competencia de este trabajo
está en la asignación, no en la variable que se consensúa, y las magnitudes
difundidas —precios, evidencia de ocupación— deben acordarse en valor
\emph{común}, no en valor opuesto. Un precio con signo contrario no tiene
lectura. El consenso bipartito no se emplea, y la razón es que la competencia ya
está internalizada en el potencial de \cref{th:potencial} en lugar de en el
signo de los enlaces.
\end{observacion}

\begin{observacion}[Retardos, pérdidas y realimentación de salida]
\label{obs:hibrido-restante}
Tres extensiones del mismo marco tocan hipótesis que este trabajo declara.
La primera, el consenso bajo retardo, cuantifica lo que \cref{pr:implementable}
supone acotado por $\Delta$. La segunda, el consenso con pérdidas de paquetes,
sostiene \cref{obs:perdidas} y admite tratamiento como sistema híbrido con
saltos en los instantes de pérdida. La tercera, la realimentación dinámica de
salida, es la más pertinente y la menos usada: los AMR no miden el estado
completo de sus vecinos, sino salidas —distancia, azimut—, y todo el aparato de
consenso de este trabajo se ha escrito sobre estados. Cerrar esa distancia exige
un observador por agente, y la elección natural es el filtro de
\cref{sec:estimacion}. Se declara como la extensión más inmediata de las tres.
\end{observacion}

\subsection{Lo que la información obsoleta cuesta en decisión}

Los resultados anteriores acotan mensajes y tiempos. Falta el paso que
convierte edad de información en error de decisión, y para un juego
fuertemente monótono ese paso es exacto \cite{mayorga2026compendio}.

\begin{teorema}[Robustez de la desigualdad variacional]
\label{th:vi-robusta}
Sea $F$ $\mu$-fuertemente monótono sobre el convexo cerrado $K$, y sean
$x^\star$ y $\hat x$ las soluciones de $\mathrm{VI}(K,F)$ y
$\mathrm{VI}(K,F+e)$ con $\lVert e(x)\rVert\le\varepsilon$ en $K$. Entonces
$\lVert\hat x-x^\star\rVert\le\varepsilon/\mu$. En particular, si el error
procede de un coste físico estimado y de precios obsoletos,
$e=\Delta c+A^{\!\top}\Delta\lambda$, entonces
$\lVert\hat x-x^\star\rVert\le(\lVert\Delta c\rVert+\lVert A\rVert_2\lVert\Delta\lambda\rVert)/\mu$.
\end{teorema}

\begin{proof}
Las dos desigualdades variacionales evaluadas cruzadamente dan
$(F(\hat x)-F(x^\star))^{\!\top}(\hat x-x^\star)\le-e(\hat x)^{\!\top}(\hat x-x^\star)$;
monotonía fuerte a la izquierda y Cauchy--Schwarz a la derecha. En cincuenta
perturbaciones aleatorias de un juego cuadrático en caja, el cociente
$\lVert\hat x-x^\star\rVert/(\varepsilon/\mu)$ máximo es $0{,}80$.
\end{proof}

\begin{proposicion}[Contracción del paso proyectado y número de condición estratégico]
\label{pr:contraccion-proyectada}
Si además $\mu I\preceq\nabla F\preceq L_FI$ en $K$, la actualización
$T(x)=\Pi_K(x-\alpha F(x))$ satisface
$\lVert T(x)-T(y)\rVert\le q(\alpha)\lVert x-y\rVert$ con
$q(\alpha)=\max\{|1-\alpha\mu|,|1-\alpha L_F|\}$; es contracción para
$0<\alpha<2/L_F$, y el paso $\alpha^\star=2/(L_F+\mu)$ da
$q^\star=(\kappa_F-1)/(\kappa_F+1)$ con $\kappa_F=L_F/\mu$. Para el juego de
tráfico con coste $J=c^{\!\top}x+\sum_zH_z(a_z^{\!\top}x)+\tfrac\epsilon2\lVert x\rVert^2$
y $0\le H_z''\le\bar h$, se tiene $\mu=\epsilon$ y $L_F\le\epsilon+\lVert A\rVert_2^2\bar h$.
\end{proposicion}

\begin{proof}
$F(x)-F(y)=\bar H(x-y)$ con $\bar H$ media de hessianas de espectro en
$[\mu,L_F]$, la proyección es no expansiva, y el mínimo de $q$ iguala
$|1-\alpha\mu|$ y $|1-\alpha L_F|$. La hessiana de $J$ es
$A^{\!\top}\mathrm{diag}(H_z'')A+\epsilon I$. En trescientos pares aleatorios
la razón observada es $0{,}930$ frente a la cota $0{,}948$.
\end{proof}

\begin{observacion}[La geometría entra por $A$, la congestión por $H''$]
\label{obs:condicion-estrategica}
$\kappa_F$ grande significa un paisaje rígido en unas direcciones y casi plano
en otras: a igualdad de integrador, condiciona la velocidad y la sensibilidad
numérica, y la regularización $\epsilon$ que lo acota es la misma que
\cref{obs:tikhonov-sesgo} paga en sesgo. Es el mismo compromiso que el
déficit de contractividad de \cref{sec:alcance}, ahora con constantes
explícitas.
\end{observacion}

\begin{proposicion}[Seguimiento por conectividad y error de decisión]
\label{pr:seguimiento-conectividad}
Sea $z$ el desacuerdo de un consenso dinámico,
$\dot z=-k(L_G\otimes I)z+d(t)$ con $z\perp\mathbf 1$ y
$\lVert d\rVert\le\bar d$. Entonces
$\limsup_t\lVert z(t)\rVert\le\bar d/(k\lambda_2(L_G))$, y si el error de
agregado entra linealmente al pago por una matriz $B$,
\begin{equation}
  \label{eq:conectividad-decision}
  \limsup_t\lVert\hat x-x^\star\rVert\le\frac{\lVert B\rVert_2\,\bar d}{\mu\,k\,\lambda_2(L_G)} .
\end{equation}
\end{proposicion}

\begin{proof}
Con $V=\lVert z\rVert^2/2$ y Rayleigh en $\mathbf 1^\perp$,
$\dot V\le-k\lambda_2\lVert z\rVert^2+\bar d\lVert z\rVert$, negativa fuera de
la bola; \eqref{eq:conectividad-decision} compone con \cref{th:vi-robusta}.
Sobre sesenta grafos conectados aleatorios el mayor cociente error/cota fue
$0{,}774$ \cite{mayorga2026compendio}.
\end{proof}

\begin{proposicion}[Obsolescencia con disparo por eventos]
\label{pr:obsolescencia}
Si el precio $\lambda_z$ se transmite cuando $|\lambda_z-\hat\lambda_z|\ge\varepsilon_z$
y $|\dot\lambda_z|\le M_z$, entonces entre transmisiones
$|\lambda_z-\hat\lambda_z|\le\varepsilon_z$, el tiempo entre eventos es al
menos $\varepsilon_z/M_z$, con retardo acotado $\Delta_z$ la obsolescencia es
a lo sumo $\varepsilon_z+M_z\Delta_z$, y apilando esas cotas en $\varepsilon$,
$\lVert\hat x-x^\star\rVert\le\lVert A\rVert_2\lVert\varepsilon\rVert/\mu$.
\end{proposicion}

\begin{proof}
Desde una transmisión el error parte de cero y crece a lo sumo a $M_z$; el
retardo añade $M_z\Delta_z$; la última cota es \cref{th:vi-robusta}.
\end{proof}

\begin{observacion}[Cierra el bucle de \cref{pr:miet}]
\label{obs:obsolescencia-cierra}
\Cref{pr:miet} acota el tiempo mínimo entre eventos; \cref{pr:obsolescencia}
acota lo que ese ahorro cuesta en decisión. Juntas dan el compromiso mensajes
frente a error de equilibrio con ambas constantes explícitas, y en el banco la
cota conserva el orden correcto al relajar el umbral aunque siga siendo
conservadora \cite{mayorga2026compendio}.
\end{observacion}

\begin{proposicion}[Seguimiento de un equilibrio móvil]
\label{pr:equilibrio-movil}
Si $F_t$ es uniformemente $\mu$-fuertemente monótono con cero interior
$x^\star(t)$ y $\lVert\dot x^\star(t)\rVert\le v$, entonces bajo
$\dot x=-F_t(x)$ se cumple $\limsup_t\lVert x(t)-x^\star(t)\rVert\le v/\mu$.
\end{proposicion}

\begin{proof}
Con $e=x-x^\star$ y $V=\lVert e\rVert^2/2$,
$\dot V=-e^{\!\top}(F_t(x)-F_t(x^\star))-e^{\!\top}\dot x^\star\le-\mu\lVert e\rVert^2+v\lVert e\rVert$.
\end{proof}

\begin{proposicion}[Frontera modal retardo--ganancia]
\label{pr:frontera-retardo}
Para el modo lineal de un lazo de precio con amortiguamiento $h$, ganancia
$\kappa$ y retardo $\tau$, $\ddot x=-h\dot x-\kappa^2x(t-\tau)$, la ecuación
característica es $s^2+hs+\kappa^2e^{-s\tau}=0$ y el primer cruce del eje
imaginario ocurre en
\begin{equation}
  \label{eq:frontera-retardo}
  \kappa^2=\omega\sqrt{\omega^2+h^2},\qquad
  \tau^\star=\frac{1}{\omega}\arctan\frac{h}{\omega}.
\end{equation}
\end{proposicion}

\begin{proof}
Con $s=j\omega$: $-\omega^2+\kappa^2\cos\omega\tau=0$ y
$h\omega-\kappa^2\sin\omega\tau=0$; dividir da $\tan\omega\tau=h/\omega$ y
elevar al cuadrado y sumar da $\kappa^4=\omega^4+h^2\omega^2$. Con $h=1$ y
$\kappa=1{,}5$: $\omega=1{,}3435$, $\tau^\star=0{,}4763$, y la raíz dominante
tiene parte real $-0{,}040$ en $\tau^\star-0{,}05$ y $+0{,}036$ en
$\tau^\star+0{,}05$.
\end{proof}

\begin{observacion}[La ventana de \cref{th:N-ventana}, con retardo]
\label{obs:retardo-ventana}
\Cref{th:N-ventana} da el paso admisible sin retardo; \eqref{eq:frontera-retardo}
dice cuánto retardo tolera cada modo de la red de precios a ganancia fija,
y cómo comprar margen: más amortiguamiento $h$ o menos ganancia $\kappa$. Es
la lectura escalar del intervalo de \cref{th:intervalo}.
\end{observacion}

\subsection{Del proceso microscópico de revisiones a la dinámica poblacional}

\begin{proposicion}[Límite de población para revisiones de Poisson]
\label{pr:limite-poblacion}
Sean $N$ agentes con estrategias $1,\dots,m$, distribución empírica $X^N(t)$
y tasas de revisión $i\to j$ iguales a $\rho_{ij}(x)$, acotadas y Lipschitz
en el símplex. Con $V(x)=\sum_{i,j}x_i\rho_{ij}(x)(e_j-e_i)$, si
$X^N(0)\to x_0$ en media cuadrática y $\dot x=V(x)$, $x(0)=x_0$, entonces
para todo horizonte $T$, $\sup_{t\le T}\lVert X^N(t)-x(t)\rVert\to0$ en
probabilidad, y con inicialización de error cuadrático $O(1/N)$ existe $C_T$
con $\mathbb E\sup_{t\le T}\lVert X^N(t)-x(t)\rVert^2\le C_T/N$
\cite{kurtz1970solutions,sandholm2010population}.
\end{proposicion}

\begin{proof}
Es la aproximación de Kurtz para cadenas dependientes de la densidad:
$X^N(t)=X^N(0)+\int_0^tV(X^N)\,\mathrm ds+M^N(t)$ con saltos de tamaño
$O(1/N)$ y variación cuadrática $O(T/N)$; restar la EDO, Lipschitz de $V$,
Doob y Gronwall. En el banco, Smith microscópico con relojes exponenciales
frente a su EDO dio pendiente log--log de la raíz del error cuadrático medio
final de $-0{,}503$ \cite{mayorga2026compendio}; una réplica reducida con
ocho semillas por tamaño reproduce el decaimiento de $0{,}028$ a $0{,}005$
entre $N=100$ y $N=6400$.
\end{proof}

\begin{observacion}[La EDO no divide coaliciones]
\label{obs:edo-macroscopica}
La ecuación de Smith no debe leerse como si una coalición pudiera dividirse
físicamente: su justificación es macroscópica, y la varianza entre semillas
cerca de un equilibrio hiperbólico es de orden $N^{-1/2}$, una propiedad
física del tamaño de la población y no ruido numérico. Es la respuesta
cuantitativa a la objeción de \cref{obs:analogia}: el símplex es invariante
en el límite, y a $N$ finito el error es $O(N^{-1/2})$.
\end{observacion}
